\chapter{Force, Inertia, and Momentum}
\label{chap:force-inertia-momentum}

The first chapter built the instrument and gave it its founding discipline: a record's physical content is
exactly what survives relabeling --- the count the world fixes, set against the bookkeeping that is ours. This
chapter sets that same discipline to work on a motion that changes. Force, inertia, and momentum are not
substances a body carries past the instrument's sensors; they are what the world charges when the instrument
re-files a record to keep it consistent. The re-filing is ours --- the instrument's own act of reconciliation,
bringing an old run of counts into alignment with a new rate --- but the bill is not: force, inertia, and
momentum are the costs the world exacts for it, a price paid to be redirected, a balance kept across a
boundary, a curvature the record already carries. Most of these costs are finite count obligations, the
instrument's ordinary honest floor; one --- the analytical formalism that carries the conjugate pair forward
--- is a smooth shadow, the continuum picture of a discrete update; and the last, at the chapter's end, is
something new: the first place where the honest floor is not a count that balances but a refinement the
instrument can prove cannot be made at all --- the book's first wall.

\section{The Newton Effect: mass as a reconciliation cost}
\label{se:newton}

This is the instrument's loveliest early idea, and it is worth taking slowly. To the instrument a body has no
inner substance to be pushed; there is only a record, a ledger of where the body has been counted and when.
Inertia is that ledger's resistance to a change in its own rate --- its reluctance to be re-filed at a new
cadence --- and mass is what that reluctance costs. The re-filing is the instrument's own work: when a force
acts, the instrument must \emph{reconcile} the record, bringing the established run of counts into alignment
with a new rate so that the old order and the new agree at their boundary. That reconciliation is never free,
and the persistent price of it --- how much re-filing each unit of redirection demands --- is exactly what we
call the body's mass. The act is the instrument's; the bill is the world's.

Newton's second law is the picture of this exchange, and it should be read as a picture and not as a recipe:
\begin{equation}
  F \;=\; m\,\ddot{x} \;=\; \dot{p}.
\end{equation}
The equation does not instruct the instrument to differentiate a momentum in order to manufacture a force. It
names a rate: the force is the rate at which the reconciliation's cost accrues as we redirect the record, and
the mass $m$ is the conversion between the force applied and the change in motion it buys. A large mass is not
a body holding more stuff; it is a record expensive to redirect --- one whose established order resists our
re-filing, so that the same applied force buys a smaller $\ddot{x}$ exactly when $m$ is large. Read this way,
$F = m\ddot{x}$ says that the world charges a fixed rate, set by the mass, to change what a record is doing;
the pushing is ours, the rate is the world's.

The bench makes it concrete. Set two carts on the track of the first chapter, one light and one heavy, and
give each the same steady push --- the same force, applied for the same interval. The instrument does not
weigh the carts; it watches their ledgers. The light cart's record is redirected quickly: a few
reconciliations bring its run of counts onto the new, faster cadence. The heavy cart's record resists --- the
same push must be paid out over many more reconciliations before its established order yields to the new rate,
and over the fixed interval its motion changes less. The instrument has not measured a substance inside either
cart; it has measured what the world charges to re-file each record --- reconciliations per unit of
redirection --- and that charge is the mass. The push was ours and the same for both; the price came back
different, and the difference was the world's to set.

The honest floor is a finite count obligation. The instrument does not certify a continuum of trajectories;
it certifies that, count for count, the change in motion the record shows is the one the force paid for. The
reading is that mass is a price the world charges, not a substance a body holds --- the cost, in
reconciliations, of changing what a record is doing, exacted no matter who does the pushing. It is the first
quantity the instrument reads that looks like a property of a body and turns out to be a property of a record
--- the world's standing charge for re-filing it, not a thing carried inside.

\section{The Momentum Effect: the conjugate that balances flux}
\label{se:momentum}

If mass is the price of redirection, momentum is the quantity whose books must balance whenever content
crosses a boundary. To the instrument, position and momentum are not two independent readings but a conjugate
pair --- two faces of one bookkeeping, locked together so tightly that the structure has a name,
\begin{equation}
  \{x, p\} = 1.
\end{equation}
This bracket is a description, not a calculation: it says that $x$ and $p$ are paired, that a shift in one is
registered against the other, that they are the canonical partners the instrument carries together. Which face
we call $x$ and which $p$, and which we write first, is ours --- but the bracket is \emph{antisymmetric},
$\{x,p\} = -\{p,x\}$, and \emph{that} reversing the order flips the sign is not ours at all. It is the same
refusal to commute that the ordered pair of the last chapter carried: compose the partners in one order and
then the reverse, and the two orderings do not agree but answer with opposite sign. The order is the hand's;
the sign-flip is the world's. The pairing will return as a central character of the quantum part, where the same conjugacy
sets the floor on how sharply the two can be read at once and builds the first quantized ladder; here it is
introduced in its gentlest, classical form, as the partnership that makes conservation possible.

Conservation, to the instrument, is a continuity of counts. Writing $\rho$ for a density of recorded content
and $j$ for its flux, the balance is
\begin{equation}
  \partial_t \rho + \nabla \cdot j = 0,
  \qquad\text{equivalently}\qquad
  \oint_{\partial \Omega} j \cdot dA = -\,\frac{d}{dt}\!\int_\Omega \rho\, dV .
\end{equation}
Neither form is a differential equation to be solved; both are descriptions of a ledger that does not leak.
The integral form says it most plainly: what leaves a region through its boundary is exactly what the region
loses, count for count. Momentum is the bookkeeping term we carry --- the conjugate to position --- to make
that balance keepable; but the balance itself is the world's, a closure it exacts whether or not we have a name
for the term that books it. Nothing crosses a boundary uncounted, and nothing is created or destroyed without
an entry: that is not a rule we impose but a debt the world collects.

Picture the boundary on the bench. Draw a closed line around a stretch of the track and let carts roll across
it, some entering, some leaving. The instrument need not watch the interior at all: it stands at the boundary
and keeps two tallies of its own, one for the momentum carried in across the line and one for the momentum
carried out. The continuity equation is the promise that these two tallies, set against the change in what the
region holds, always close --- that a region's total cannot move except by what crosses its edge, counted in
and counted out. Momentum is whatever makes that promise keepable: the conjugate quantity whose inflow and
outflow the instrument can balance against the region's change without ever opening the region up. The tallies
are the instrument's to keep; that they are forced to close is the world's doing --- and the same trick of
judging an interior from
its edge alone is the seed of the boundary reconciliations that will run through the quantum and relativity
parts, where a record is fixed from its frontier without the interior being scanned.

The honest floor is a finite count obligation: a flux that must balance across a boundary, count for count.
The reading is that momentum is not a thing a body ferries through space but the conjugate that keeps the
books balanced as content crosses a boundary --- the books are ours, the balance is the world's, and
conservation is the name of the second. The instrument does not enforce conservation as an external rule; it
reads it off the requirement that a boundary's accounts close.

\section{The Lagrange--Hamilton Effect: the path of stationary action}
\label{se:lagrange-hamilton}

The conjugate pair the last effect named has a dynamics, and developing it is the single most valuable thing
this chapter does, because the structure laid down here is the classical precursor of the quantum part's
deepest machinery. Begin with the Lagrangian, the difference between a record's kinetic and potential
bookkeeping,
\begin{equation}
  L = T - V,
\end{equation}
and its accumulation along a candidate motion, the action,
\begin{equation}
  S = \int L\, dt .
\end{equation}
Both are ours to write: the split into kinetic and potential, and the very list of candidate paths a record
might draw between two fixed endpoints, are a frame the instrument lays over the motion, not a thing the motion
announces. What is not ours is which path is singled out. Among all the candidates we propose, the one the
world realizes is the one for which the action is stationary --- the path whose neighbors, to first order, cost
the same,
\begin{equation}
  \delta S = 0 .
\end{equation}
This is the least-principle the book returns to two chapters on, named here at its first appearance: the
realized motion is not pushed step by step but selected whole, as the stationary path of an accumulated cost.
We supply the candidates and the cost we score them by; the world supplies the one that stands still under
variation. The instrument reads it as an extremal count --- the motion whose total bookkeeping no small
re-routing improves.

The same content has a second form, dual to the first and closer to how the instrument actually steps. Where
the Lagrangian is the difference of the two energies, the Hamiltonian is their sum,
\begin{equation}
  H = T + V,
\end{equation}
the total energy written as a function of the conjugate pair $(q, p)$ --- position and its partner momentum.
And the Hamiltonian generates the instrument's update rule. Given the record's current $(q, p)$, the next is
fixed by Hamilton's equations,
\begin{equation}
  \dot q = \frac{\partial H}{\partial p}, \qquad \dot p = -\frac{\partial H}{\partial q},
\end{equation}
a pair that says, descriptively, how each member of the conjugate pair is re-filed in terms of the other:
position advances at a rate set by how the energy varies with momentum, momentum at a rate set, with a sign,
by how the energy varies with position. The smooth equations are our writing of it; the step itself --- one
count of the conjugate pair carried to the next --- is what the instrument actually does, and the world is what
makes the carry come out one way and not another. This is the instrument's step, not a formula to evaluate but
the rule by which the record is advanced.

The pairing itself governs how every quantity rides this flow. For any reading $f$ built on the conjugate
pair, its rate of change along the motion is its bracket with the Hamiltonian,
\begin{equation}
  \dot f = \{f, H\},
\end{equation}
the bracket measuring how $f$ co-varies with the conjugate structure as the energy drives it forward. And here
the forward-link must be stated plainly, because it pays off the whole quantum part. The conjugate bracket the
last effect introduced, $\{q, p\} = 1$, is the classical shadow of the relation the quantum chapter is built
on, the canonical commutator
\begin{equation}
  [x, p] = i\hbar .
\end{equation}
The two are the same conjugate pairing at two scales --- the bracket the smooth, classical version, the
commutator the discrete, quantized one, with $\hbar$ the size of the rung between them. The names are ours; the
pairing they both fix, and its refusal to commute, is the world's at either scale. Everything the quantum part
raises on the commutator --- the uncertainty floor, the oscillator's ladder, the reversible evolution --- is
already foreshadowed here, in the bracket that pairs position with momentum.

The honest floor is a smooth-shadow analogy: the discrete update of the conjugate pair, count by count, is the
model; the continuum formalism --- the action integral, Hamilton's equations, the smooth flow --- is its
shadow, ours to draw and named as one. The reading is that the analytical formalism is the instrument's
stepping rule written smooth: we write the Lagrangian and the flow; the world selects the stationary-action
path; and the Hamiltonian we wrote carries the conjugate pair forward only because the world lets the carry
close. The bracket $\{q, p\} = 1$ holds the pairing fixed --- the same pairing that, sharpened by $\hbar$,
becomes the commutator the quantum part is built on.

\section{The Acceleration Effect: curvature, not force}
\label{se:acceleration}

Acceleration is usually introduced as a force divided by a mass, but the instrument meets it earlier than any
force, as something purely geometric: the curvature of the trace the record draws. It is the second difference
of position --- how far the trajectory bends away from a straight continuation ---
\begin{equation}
  a \;=\; \ddot{x} \;=\; \lim_{\Delta t \to 0}\frac{x_{n+1} - 2x_n + x_{n-1}}{\Delta t^2}.
\end{equation}
The quantity that matters here is the bare second difference $x_{n+1} - 2x_n + x_{n-1}$, and it should be read
as a shape, not a derivative to be evaluated. It is the gap between where the record actually lands and where a
straight-line continuation of the previous two counts would have put it. The straight continuation is ours: it
is the standard of \emph{no bend} we lay down, the frame against which a bend is even measured. What fills the
gap is not. A body gliding in a straight line has $a = 0$ identically, however fast it moves, and no change of
our frame puts a bend into it; a body that bends carries that bend in every frame at once. Acceleration is the
instrument's count of the ways a straight-line fit to the record fails --- informational curvature, the bend in
a ledger, computed from the record alone, with no force named and none needed. The standard is the hand's; the
departure from it is the world's.

This descriptive reading is the seed of one of the book's later pictures. In the relativity part the same
curvature, no longer of a trace in space but of the record's own order, will become gravity: a body in free
fall draws the straightest line its curved bookkeeping allows, and what we call the gravitational force is a
name we assign to a curvature the record already carries, no frame of ours able to straighten it away. Here, in
flat bookkeeping, the point is simpler and prior --- the bend is the world's, the force is our word for it, and
the bend comes first.

The honest floor is a finite count obligation: the second difference is a count of fit-failures, secured from
the record without invoking dynamics. The reading is that acceleration is curvature in the record --- what
bends, not what pushes. The straight line we fit is ours; the bend that defeats the fit is the world's; and a
force, when it is later introduced, is the cause we assign to a curvature the instrument has already read off
the record.

\section{The da Vinci--Coulomb Effect: the bound that is not a procedure}
\label{se:davinci-coulomb}

Here the book meets its first impossibility, and the character of the instrument's honest floor changes.
Static friction holds a body still until an applied force crosses a threshold, an effect Leonardo noted and
Amontons and Coulomb made law: a body sticks while the tangential force satisfies
\begin{equation}
  |F| \le \mu_s\, |N|,
\end{equation}
and slips when it exceeds $\mu_s |N|$, with a coefficient $\mu_s$ that depends on the load but, famously, not
on the apparent area of contact. The bound is real and correct. But notice what kind of statement it is. It is
a \emph{bound}, not a \emph{procedure}. The inequality describes a wall --- it says where slip becomes
possible --- but it does not hand over $\mu_s$ from any single record, and below the threshold it does not
predict when, within the stuck regime, anything happens, because nothing does. The coefficient is the world's,
recoverable only as an invariant forced by many trials and never read off one; the wish to take it from a
single look is ours, and the world declines it. The threshold is well defined and yet not executable from one
observation: the instrument can confirm the wall is there without being able to locate it from a single look.

And below the threshold, slip is not merely absent but \emph{inadmissible}. There is no refinement of the
record --- no finer resolution we might bring, no closer look --- that yields a permitted slip while the bound
holds. The finer resolutions are ours to attempt; the refusal is the world's, and no labeling of ours softens
it. This is the book's first no-go, and the shift deserves marking, because it is the second of the
instrument's four honest floors to appear. The floors of the previous three effects were finite count obligations:
quantities the instrument requires to balance. The floor here is an impossibility --- a refinement the
instrument can prove inadmissible, a slip that cannot be made below threshold no matter how finely the record
is resolved. The ceiling is not a count but a wall.

The reading is that friction's threshold is a genuine feature of the world, and the law is honest precisely
about what it cannot do. It bounds the onset of slip without pretending to predict it; it marks that the wall
is there without locating, from one look, exactly where. The forecast the hand wants, the world withholds; the
bound the hand did not ask for, the world insists on. That candor --- a constraint that declares itself a
constraint and not a forecast --- is the instrument's first encounter with a kind of honesty it will meet
again at every wall in the book: crossed polarizers that extinguish a beam, a domain seam that cannot be
misdrawn, a closed loop in time the count's own rising forbids. The no-go is rarer than the count, but it is
no less exact, and the instrument names it as carefully as it names a balance.

\bigskip

With force, inertia, and momentum the instrument has begun to act, and already three of its four honest floors
are in play. Mass is the price the world charges a record to be redirected; momentum is the conjugate that
balances its books across a boundary; acceleration is the curvature it reads before any force is named --- three
finite count obligations, three balances the world exacts and the instrument secures count for count. Between
them the analytical formalism --- stationary action, the Hamiltonian update, the conjugate bracket --- is the
smooth shadow we draw over the instrument's discrete step, and it carries the conjugate pairing, ours to write
and the world's to keep antisymmetric, that the quantum part will sharpen into a commutator. And friction is
the first wall: a bound the instrument records and respects but cannot resolve below, the world's \emph{no} to
a procedure the hand would have preferred.
The next chapter carries the instrument from straight-line motion to motion that closes on itself, where the
record becomes a loop --- and where the bend it learned to read here returns as the residue a closed loop
carries.

\bigskip
\begin{center}
\rule{0.4\textwidth}{0.4pt}
\end{center}
\bigskip

\noindent A charge, once entered, must be accounted for --- and the strongest evidence a case can carry is a
quantity that cannot quietly go missing. What the record conserves is exactly that: a tally the world will not
let vanish between the moment it is written and the moment it is answered for. Nothing crosses the boundary of
the account uncounted; what enters one side leaves the other, and the books are made to close. The instrument
keeps the tally; that it must balance --- that no entry disappears without a matching one --- is the world's
insistence, not the clerk's. A charge that balances is a charge still standing when it is called.
